\documentclass[../../main.tex]{subfiles}
\graphicspath{{\subfix{../../image/}}}

\begin{document}

\section{概述}\label{section:ch01-overview}

本教材分为三个大部分。

\textbf{第一部分：解的表示公式.} 在此我们挑出那些重要的偏微分方程，即在某些情形下可以得到解的显式或近乎显式公式的方程。论述的大体进程是：从某几类线性方程的直接公式，过渡到对各类非线性 PDE 而言不那么具体的、某种意义上的表示公式。

\cref{chapter:ch02}详细研究四个可精确求解的偏微分方程：线性输运方程、Laplace方程、热方程和波动方程。这些 PDE 作为后文更复杂方程的原型，至少在不存在边界几何带来麻烦的区域的情形下，可以直接算出解。显式公式还辅以各种间接但简单而动人的``能量''型论证，它们为\cref{chapter:ch06-二阶椭圆型方程,chapter:ch07-线性发展方程}及以后的发展提供了动机。

\cref{chapter:ch03}延续寻找显式公式的主题，不过现在是对一般的一阶非线性 PDE。关键的洞察是：这类 PDE 至少在局部可以转化为常微分方程（ODE）组，即特征方程组。我们约定：一旦问题``仅仅''是积分一个 ODE 组的问题，它原则上就已求解，有时还相当地显式。正文中推导特征方程的过程非常简单，不需要任何几何洞见。事实上，推导特征方程是如此容易，以至于先处理拟线性情形并没有什么实际意义。

我们还引入了 Hamilton--Jacobi 方程的 Hopf--Lax 公式（\cref{section:ch03-hj}）以及标量守恒律的 Lax--Oleinik 公式（\cref{section:ch03-conservation}）。（这里用到一些测度论知识会有帮助，但并非必不可少。）这些小节使读者尽早熟悉这些重要非线性 PDE 的整体理论，从而为后文\cref{chapter:ch10-HamiltonJacobi方程,chapter:ch11-守恒律方程组}提供动机。

\cref{chapter:ch04}是一个技巧汇编，用于显式（或某种意义上显式）地求解各类线性与非线性偏微分方程，读者只需选读其中感兴趣的部分即可。不过，关于 Fourier 变换的那一节是必不可少的。Cauchy--Kovalevskaya 定理出现在最后。尽管它基本上是本学科中唯一的一般存在性定理，因而在逻辑上或许应视为中心内容，但在实践中这些幂级数方法并不那么盛行。

\textbf{第二部分：线性偏微分方程理论.} 接下来我们放弃对显式公式的寻求，转而依靠泛函分析和相对简单的``能量''估计，来证明各类线性 PDE 弱解的存在性。我们还研究这些解的唯一性与正则性，并推导出其他各种性质。

\cref{chapter:ch05-Sobolev空间}介绍Sobolev空间，这是用能量方法研究许多线性与非线性偏微分方程的恰当框架。这是一章硬内容，其真正价值要到后面才显现出来，并且需要一些 Lebesgue 测度论的基本知识。不过要求其实并不那么高，\cref{chapter:appE-测度论} 中的回顾应当足够。在我看来，只考察指数 $p=2$ 的Sobolev空间并无特别好处，而拘泥于此反而会遮蔽两个核心不等式，即 Gagliardo--Nirenberg--Sobolev 不等式（\cref{subsection:ch05-GNS}）与 Morrey 不等式（\cref{subsection:ch05-Morrey}）。

在\cref{chapter:ch06-二阶椭圆型方程}中，我们把关于Laplace方程的知识大幅推广到其他二阶椭圆型方程。在此我们较为完整地讨论了解的存在性、唯一性与正则性理论，包括极值原理，并对特征值的研究作了相当的介绍，其中讨论了非自伴算子的主特征值。

\cref{chapter:ch07-线性发展方程}把能量方法推广到刻画时间演化的各类线性偏微分方程。我们把先前对热方程的考察推广到一般的二阶抛物型 PDE，把对波动方程的考察推广到一般的二阶双曲型 PDE。我们还研究线性一阶双曲组，旨在为\cref{chapter:ch11-守恒律方程组}关于非线性守恒律组的发展提供动机。最后\cref{section:ch07-4-半群理论}给出建立解的另一种泛函分析方法，即半群方法。

（在这一大部头的第二部分中，略去了分布理论与位势理论的任何讨论。这些都是重要课题，但对我们的目的而言似乎可以省略，即便在这样一部篇幅的书中也是如此。这些省略并不会太妨碍我们，反而为更多的非线性理论留出了篇幅。）

\textbf{第三部分：非线性偏微分方程理论.} 这一部分就非线性 PDE 而言与第二部分相平行，但其方法远不如第二部分统一，因为各类非线性性必须以相当不同的方式处理。

\cref{chapter:ch08-变分法}以对变分法的广泛讨论开始对非线性偏微分方程的一般研究。在此我们仔细推导了用直接方法证明极小化子存在性的过程，并讨论了各种变分方程组与约束问题，以及极小极大方法。变分理论是非线性 PDE 诸方法中最有用、最易上手的一种，因此这一章是基础性的。

\cref{chapter:ch09-非变分技术}与早先的\cref{chapter:ch04}颇为相似，汇集了用于非线性椭圆型与抛物型偏微分方程的其他各种技巧。我们在此遇到单调性方法与不动点方法，以及其他各种工具，其中大多涉及极值原理。我们还研究非线性半群理论的某些漂亮内容，以补充\cref{chapter:ch07-线性发展方程}的线性半群理论。

\cref{chapter:ch10-HamiltonJacobi方程}介绍 Hamilton--Jacobi PDE 的现代理论，特别是``粘性解''的概念。我们还通过动态规划，讨论其与 ODE 最优控制的联系。

\cref{chapter:ch11-守恒律方程组}从\cref{chapter:ch03}接续守恒律的讨论，不过现在是守恒律方程组。与\cref{chapter:ch05-Sobolev空间,chapter:ch06-二阶椭圆型方程,chapter:ch07-线性发展方程,chapter:ch08-变分法,chapter:ch09-非变分技术}中以Sobolev空间为恰当抽象框架的一般理论发展不同，在此我们不得不直接使用线性代数与微积分计算。我们特别关注Riemann问题的求解与熵准则。

\cref{chapter:ch12-非线性波动方程}是对非线性波动方程的介绍，为本版新增。我们对某几类拟线性波动方程给出长时间与短时间存在性定理，并深入考察半线性波动方程，尤其是三维空间中次临界与临界幂非线性的情形。为补充这些存在性定理，最后一节指出了确保解不存在的各种准则。

\cref{chapter:appA-记号}--\cref{chapter:appE-测度论} 为读者方便起见，提供了关于不等式、线性泛函分析、测度论等的一些背景材料，并附有部分证明。

参考文献是一份经过更新的、广泛的有趣 PDE 书目，可供进一步查阅。由于这是一本教材而非专著，我大多没有去追溯并记录我们将遇到的无数思想与方法的原始出处。偏微分方程的数学文献确实浩如烟海，但参考文献中列出的书籍至少应能为查找原始文献提供一个起点。（正文中各处也引用了若干精选的研究论文。）

\end{document}